
%قبل از شروع مسئله ترجیحاً چیزی ننویسید. اما اگر خواستید بنویسید، در قالب لیستی (فقرات و شمارش) بنویسید.
%برای اضافه کردن نمره، بعدی از شروع‌مسئله، [x] را اضافه کنید و به جای x نمره مورد نظر را قرار دهید. برای زیرمسئله‌ها هم به همین صورت است.
%\شروع‌مسئله[10]
\شروع‌مسئله [18]
\\
$n$ درخت باینری داریم که هر کدام $k$ راس دارند.عملیات تعویض یعنی دو زیر درخت یک راس از یک درخت را با یکدیگر عوض کنیم. دو درخت قابل تبدیل شدن به یکدیگرند اگر و تنها اگر با تعدادی عملیات تعویض بتوان از یک درخت ، درخت دیگر را ساخت. الگوریتمی کارا از $O(nk)$ برای یافتن "تعداد" جفت درخت هایی که قابل تبدیل به یکدیگرند ارائه دهید.
%\متن‌سیاه{سوال عملی اول:}\\
%\متن‌سیاه{دراپ‌تست}\\

%\متن‌سیاه{ورودی:}\\
%در خط اول ورودی، ابتدا k و سپس m داده می‌شود.\\
%\متن‌سیاه{خروجی:}\\
%در تنها سطر خروجی، تعداد آیفون‌های مورد نیاز می‌آید. اگر در هیچ صورتی چنین کاری ممکن نبود ۱- چاپ می‌شود.\\
%\متن‌سیاه{محدودیت‌ها:}\\
%$1\leq k\leq 10^4$

%$m\leq10^9$
%\\
%\textbf{\متن‌سیاه{ورودی و خروجی نمونه}\\
%% تست اول
%\begin{latin}
% \begin{tabularx}{\textwidth}{|X | X|}
% \hline
% \textbf{stdin} &
% \textbf{stdout}\\
% \hline
% 100 100&
% 1\\
% \hline
% 
% \end{tabularx}
%\end{latin}
%
%% تست دوم
%
%\begin{latin}
% \begin{tabularx}{\textwidth}{|X | X|}
% \hline
% \textbf{stdin} &
% \textbf{stdout}\\
% \hline
% 340 20&
% 3\\
% \hline
% 
% \end{tabularx}
%\end{latin}}
% صورت اصلی مسئله را در این‌جا بنویسید.



\solution{
	درخت ها را به صورت زیر در هم سازی می‌کنیم:
	
	فرض کنید $H(v)$ درهم سازی درختی ریشه‌دار از $v$ باشد و $v$ دو تا بچه $L$ و $R$ دارد.
	کاری که می‌کنیم دو تا بچه‌را بر اساس تابع درهم سازی‌شان مرتب می‌کنیم، یعنی برای مثال اگر زیر درخت سمت چپ تابع درهم سازی اش از سمت راستی بزرگتر بود آن‌ها را جابه‌جا می‌کنیم. فرض کنید $L'$ و $R'$ ترتیب مرتب‌شده این‌دو باشد.
	$$H(v) = asciicode('(') \times mod^(size(L') + size(R')) + H(L') \times mod^(size(R) + 1) + asciicode(')')$$
	
	مقدار $H(v)$ اگر رأس $v$ وجود نداشته‌باشد برابر صفر است، همچنین $size(v)$ بیانگر اندازه زیر درخت $v$ است. توجه کنید تابع درهم سازی بالا به نوعی درخت را متناظر یک پرانتز‌گذاری می‌کند و پرانتز‌گذاری مورد نظر که یک رشته‌است از طریق روش درهم‌سازی معمول درهم سازی می‌شود. به سادگی دیده می‌شود دو درخت قابل تبدیل به هم دقیقا  یک مقدار می‌گیرند.

توجه کنید پس از درهم‌سازی کافیست یک آرایه مثل
$cnt_h$
داشته‌باشیم که تعداد درخت‌هایی را بشمارد که مقدار تابع درهم سازی آن برابر $h$  است. سپس جواب مسئله برابر مقدار زیر خواهد بود:
$$\sum_{h \in R_h} \binom{cnt_h}{2}$$
محاسبه مقدار درهم سازی هر درخت با اردر تعداد رئوس درخت یعنی $O(n)$ امکان‌پذیر است پس محاسبه مقدار بالا برابر $O(nk + mod)$ است که اگر اندازه $mod$ را مقداری منطقی انتخاب کنیم و یا درهم سازی مناسبی انتخاب کنیم الگوریتمی کارا‌ است.

توجه کنید می‌توانید $mod$ را بزرگ تعریف کنید و مقادیر مختلف درهم‌سازی را در یک درخت جست‌وجوی دودویی بریزید و مقدار 
$cnt_h$
را با
$O(log_2n)$
بشمارید. این هم مورد قبول است. 
هر تابع درهم سازی مناسب دیگر که جواب‌گو باشد برای درخت‌ها نیز جوابگو خواهد بود.

}

%%%%%%%%%%%%%%%%%%%%%%%%%%%%%%%%%%%%%%%%%%%%%%%%%%%%%
% برخی دستورات مفید برای تمرین‌ها. %
%%%%%%%%%%%%%%%%%%%%%%%%%%%%%%%%%%%%%%%%%%%%%%%%%%%%%

% فرمول‌های نمونه:
%$f(x) = \Theta(x)$
%$T(x, y) \leq \log(x) + 2^y + \binom{2n}{n}$
% {\tt code\_function}


% لیست موارد بدون شماره‌گذاری:
%\شروع{فقرات}
% \فقره مورد اول 
% \شروع{فقرات}
% \فقره مورد اول داخلی 
% \شروع{فقرات}
% \فقره اولین مورد خیلی داخلی 
% \فقره دومین مورد خیلی داخلی 
% \پایان{فقرات}
% \فقره مورد دوم داخلی 
% \پایان{فقرات}
% \فقره مورد دوم 
%\پایان{فقرات}

% لیست موارد شماره‌گذاری شده:
%\شروع{شمارش}
% \فقره یک مورد 
% \فقره یک مورد دیگر 
% \فقره باز هم یک مورد دیگر
%\پایان{شمارش}

% برای نوشتن یک فرمول چند خطی که یک شماره بخوره. برای حالت‌های دیگر به سایت زیر مراجعه کنید:
% https://www.overleaf.com/learn/latex/Aligning_equations_with_amsmath
%\begin{align} 
% \begin{split}
% (x+y)^3 &= (x+y)^2(x+y) \\
% &= (x^2+2xy+y^2)(x+y) \\
% &= (x^3+2x^2y+xy^2) + (x^2y+2xy^2+y^3) \\
% &= x^3+3x^2y+3xy^2+y^3
% \end{split} 
%\end{align}

% نوشتن ماتریس و آرایه دوبعدی
%\begin{align}
% A = 
% \begin{bmatrix}
% A_{11} & A_{21} \\
% A_{21} & A_{22}
% \end{bmatrix}
%\end{align}

% اضافه کردن تصویر با شرح. برای اضافه کردن تصویر بدون شرح، فقط دستور درج‌تصویر را استفاده کنید.
%\شروع{شکل}[h!]
%\تنظیم‌ازوسط
% \درج‌تصویر[پهنا=0.5\پهنای‌متن]{img.jpg}
%\شرح{شرح تصویر}
%\برچسب{تصویر:۱}
%\پایان{شکل}
%\پایان{زیرمسائل}

%کد زیر را ببینید.
%%\lstinputlisting{source_filename.py}
%\begin{latin}
%\begin{lstlisting}
%# Import the modules
%import sys
%import random
%
%ans = True
%
%while ans:
%question = raw_input("Ask the magic 8 ball a question: (press enter to quit) ")
%
%answers = random.randint(1,8)
%
%if question == "":
%sys.exit()
%
%elif answers == 1:
%print "It is certain"
%
%elif answers == 2:
%print "Outlook good"
%
%elif answers == 3:
%print "You may rely on it"
%
%elif answers == 4:
%print "Ask again later"
%
%elif answers == 5:
%print "Concentrate and ask again"
%
%elif answers == 6:
%print "Reply hazy, try again"
%
%elif answers == 7:
%print "My reply is no"
%
%elif answers == 8:
%print "My sources say no"
%\end{lstlisting}
%\end{latin}